% ==============================================================
% DISCUSIÓN Y CONCLUSIONES — RESULTADOS DE AMENAZA Y RIESGO
% ==============================================================
\chapter{Discusión y conclusiones}
\label{ch:discusion_conclusiones}

\section{Amenaza sísmica}
\label{sec:disc_amenaza}

\subsection{Amenaza determinista (DSHA)}
\label{subsec:disc_dsha}

\subsubsection{Espectros en sitios y contraste con la referencia normativa}
\label{subsubsec:disc_dsha_espectros}

Los espectros DSHA en cinco sitios (Figura~\ref{fig:dsha_espectros_5sitios}) muestran un patrón consistente con un control por cercanía y contenido de frecuencia: en los Sitios 1, 2, 3 y 5 las mayores ordenadas se concentran en periodos cortos ($\sim0.1$--$0.4$~s), mientras que el Sitio 4 presenta un comportamiento distinto en la referencia normativa, con un espectro elástico NCh433 más elevado en periodos cortos a intermedios. Esta diferencia no es un detalle menor, porque condiciona qué tan informativo resulta el contraste ``escenario versus norma'' y, por tanto, cómo se interpreta la severidad relativa del escenario determinista en un emplazamiento específico.

La Figura~\ref{fig:dsha_diff_espectros_nch433} muestra que en los Sitios 1 y 2 existen excedencias positivas en periodos cortos para varias curvas. En términos ingenieriles, esto es incisivo: en esos emplazamientos, el DSHA identifica escenarios plausibles con demanda sobre la referencia normativa en el rango donde típicamente se controla el diseño de estructuras rígidas o de baja altura. En cambio, en el Sitio 4 predominan diferencias negativas de mayor magnitud, indicando que allí la norma impone una demanda más alta en periodos cortos a intermedios que la derivada de los escenarios analizados. La lectura que se desprende es que el contraste con NCh433 no entrega un mensaje global único; revela, más bien, una interacción entre emplazamiento y forma espectral que puede invertir la relación ``más demandante'' según el sitio.

El resumen de PGA y $SA(1.0)$ (Figura~\ref{fig:dsha_resumen_pga_sa10_sitios}) refuerza el punto clave: la jerarquía espacial es más nítida en $SA(1.0)$ que en PGA. Esto implica que el orden relativo entre ubicaciones no está fijado por un único indicador de alta frecuencia, sino que puede cambiar al pasar a periodos intermedios. En la práctica, este resultado obliga a ser cuidadoso al resumir escenarios cercanos mediante PGA: diferencias relativamente pequeñas en PGA pueden traducirse en separaciones más amplias en $SA(1.0)$ y, por ende, en demandas más contrastadas para tipologías cuya respuesta está gobernada por periodos del orden de 1~s.

\subsubsection{Mapas DSHA para escenarios principales}
\label{subsubsec:disc_dsha_mapas_principales}

Los mapas \pFifty{} de PGA y $SA(1.0)$ para los escenarios principales (Figuras~\ref{fig:dsha_map_pga_all} y \ref{fig:dsha_map_sa10_all}) muestran tres huellas espaciales con rasgos robustos. Interplaca presenta un campo suave, con variación espacial moderada, compatible con una fuente extensa y distante cuyo efecto se distribuye regionalmente. Intraplaca exhibe valores relativamente altos y más homogéneos a escala de cuenca, coherente con un escenario profundo y de gran magnitud, donde el \textit{footprint} domina sobre contrastes locales. En contraste, FSR produce un gradiente longitudinal pronunciado y máximos hacia el margen oriental del dominio, alineados con la traza, característica inequívoca de control por distancia a una fuente cortical cercana.

Este contraste es decisivo para la lectura territorial: incluso si las fuentes de subducción resultan dominantes al integrar recurrencias en un marco probabilístico, el DSHA demuestra que la FSR organiza la demanda máxima plausible de manera concentrada en el piedemonte. En otras palabras, la FSR no compite ``en promedio regional'' con subducción; compite localmente donde la distancia mínima se reduce, y allí el gradiente es lo suficientemente fuerte como para reordenar prioridades espaciales.

\subsubsection{Sensibilidades geométricas: intraplaca versus FSR}
\label{subsubsec:disc_dsha_sensibilidades}

La sensibilidad intraplaca a la rotación del dip (Figura~\ref{fig:dsha_map_pga_intra_profundidad}) muestra patrones regionales similares con variaciones localizadas. El mensaje es concreto: para el dominio analizado, la incertidumbre geométrica intraplaca redistribuye intensidades en forma secundaria respecto del control dominante de una fuente amplia y profunda. Esto no implica irrelevancia, pero sí sugiere que, en la escala de cuenca, el error inducido por estas variaciones es moderado frente a la magnitud de la señal intraplaca misma.

En FSR ocurre lo contrario. La sensibilidad por traza y dip (Figuras~\ref{fig:dsha_map_pga_fsr_dip_traza} y \ref{fig:dsha_map_sa10_fsr_dip_traza}) produce cambios más notorios, concentrados especialmente hacia el margen oriental. Esto es esperable, pero aquí se cuantifica espacialmente: en una fuente somera, cambiar traza o inclinación modifica directamente el campo de distancias cercanas y desplaza máximos en franjas estrechas. El resultado es que la incertidumbre geométrica de FSR tiene implicancias operativas inmediatas para lectura local de demanda (y posteriormente de pérdidas), mientras que en intraplaca su efecto tiende a diluirse en un patrón regionalmente coherente.

\subsubsection{Percentiles GMF y perfiles longitudinales}
\label{subsubsec:disc_dsha_percentiles_perfiles}

Los percentiles \pSixteen{}, \pFifty{} y \pEightyFour{} del conjunto GMF en FSR (Figura~\ref{fig:dsha_map_pga_fsr_percentiles}) muestran una expansión sistemática de zonas de alto PGA hacia el este al pasar de \pSixteen{} a \pEightyFour{}. La señal no es sólo ``más alto en \pEightyFour{}'': la geometría espacial del campo intenso se expande donde la cercanía a la traza vuelve a la predicción más sensible. Esta evidencia es relevante para riesgo porque incrementos en percentiles de intensidad en la franja de mayor demanda son precisamente los que más pueden activar no linealidades de daño y pérdida.

Las diferencias entre familias de fuentes (Figura~\ref{fig:dsha_map_pga_diff_subduccion}) muestran cambios de signo longitudinales, confirmando una transición de control: subducción tiende a dominar en sectores occidentales donde FSR es más distante, mientras que FSR domina en la franja oriental donde la distancia disminuye. Los perfiles Este--Oeste (Figuras~\ref{fig:dsha_perfiles_ew_pga} y \ref{fig:dsha_perfiles_ew_sa10}) y sus bandas \pSixteen{}--\pEightyFour{} (Figuras~\ref{fig:dsha_perfiles_ew_pga_bandas_lat33.4}--\ref{fig:dsha_perfiles_ew_sa10_bandas_lat33.6}) consolidan esta interpretación al mostrar que interplaca e intraplaca varían suavemente con la longitud, mientras que FSR crece de forma marcada hacia el este y, además, ensancha su banda de incertidumbre en esa dirección. Esto implica que el extremo oriental combina mayor demanda central y mayor dispersión, exactamente la combinación que tiende a amplificar el riesgo local.

\subsection{Amenaza probabilística (PSHA)}
\label{subsec:disc_psha}

\subsubsection{Curvas PSHA: dominancia de subducción y rol acotado de FSR}
\label{subsubsec:disc_psha_curvas}

Las curvas PSHA (Figuras~\ref{fig:psha_curvas_pga_poisson}--\ref{fig:psha_curvas_sa10_bpt}) muestran que, para el rango de PoE reportado, la curva total está controlada por el envolvente de subducción, mientras que FSR aparece como contribución menor. Este resultado es consistente con la interpretación probabilística de $\lambda(IM>im)$ (Ec.~\ref{eq:hazard_curve_def}) y su traducción a PoE en un horizonte finito (Ec.~\ref{eq:poe_from_lambda}): al incorporar tasas y una familia amplia de escenarios en subducción, la contribución agregada de interplaca e intraplaca tiende a dominar, incluso si una fuente cortical puede ser muy intensa localmente en escenarios específicos.

El contraste entre PGA y $SA(1.0)$ es particularmente informativo: $SA(1.0)$ muestra una dominancia interplaca más marcada que PGA en el mismo conjunto de sitios, lo que anticipa que los mecanismos que controlan amenaza a periodos intermedios son más dependientes de eventos grandes y relativamente distantes que de contribuciones corticales cercanas.

\subsubsection{Desagregación en PGA: reparto interplaca/intraplaca y señal de FSR}
\label{subsubsec:disc_psha_disagg_pga}

La desagregación por fuente para PGA (Figuras~\ref{fig:psha_disagg_pga_poisson} y \ref{fig:psha_disagg_pga_bpt}) muestra que, bajo Poisson, el reparto interplaca/intraplaca es efectivamente competitivo en los Sitios 1--3, con FSR en niveles bajos pero no nulos. Para PoE 2\% en 50 años se reportan, por ejemplo, contribuciones FSR de 1--3\% en Sitios 1--3, con el resto repartido casi por mitades entre interplaca e intraplaca; en Sitios 4 y 5, FSR es nula y el reparto se concentra completamente en subducción. Esta asimetría espacial es incisiva: la amenaza cortical no sólo es ``pequeña en promedio'', sino que además se apaga en sitios donde la distancia efectiva hace irrelevante su contribución en el rango de PoE analizado.

Bajo BPT, la contribución FSR en PGA aumenta en Sitios 1--3 (por ejemplo, a 3--6\% según sitio y PoE), mientras Sitios 4 y 5 se mantienen esencialmente sin contribución cortical. Esto implica que la dependencia temporal no redistribuye el control hacia FSR de manera uniforme; realza la fuente cortical precisamente donde ya tenía presencia (sitios más cercanos o más sensibles a ella), reforzando una lectura localizada del impacto de modelos temporales.

\subsubsection{\texorpdfstring{Desagregación en $SA(1.0)$: dominancia interplaca y límites del aporte cortical}{Desagregacion en SA(1.0): dominancia interplaca y limites del aporte cortical}}
\label{subsubsec:disc_psha_disagg_sa10}

La desagregación para $SA(1.0)$ (Figuras~\ref{fig:psha_disagg_sa10_poisson} y \ref{fig:psha_disagg_sa10_bpt}) es aún más concluyente. Bajo Poisson y PoE 2\% en 50 años, la contribución interplaca es del orden 79\%--89\% en los cinco sitios, con intraplaca típicamente en 10\%--16\% y FSR entre 0\% y 6\%, con valores nulos o cercanos a cero en Sitios 4 y 5. Este resultado fija un límite interpretativo: en el rango de amenaza probabilística representado por $SA(1.0)$, la FSR puede ser relevante localmente, pero no desplaza el control interplaca, y su influencia queda confinada a sitios específicos.

Bajo BPT, la contribución FSR en $SA(1.0)$ aumenta de forma apreciable en Sitios 1 y 2, alcanzando 13\% en PoE 2\% en 50 años, mientras permanece menor en Sitios 3--5. El mensaje no es que ``FSR pasa a dominar'', sino que su aporte relativo puede ser material en periodos intermedios en sitios cercanos, aun cuando interplaca siga siendo el componente mayoritario. Por lo tanto, el modelo temporal no cambia el diagnóstico de dominancia en $SA(1.0)$, pero sí amplía el rango plausible de participación cortical en la amenaza en sectores específicos.

\subsubsection{Mapas PSHA: magnitud del efecto de incluir FSR y su localización}
\label{subsubsec:disc_psha_mapas}

Los mapas PSHA de PGA para PoE 2\% y 10\% en 50 años (Figuras~\ref{fig:psha_map_pga_poe02} y \ref{fig:psha_map_pga_poe10}) muestran que la estructura espacial global es estable entre Poisson completo, Poisson sin FSR y BPT completo. Las diferencias por inclusión de FSR (Figuras~\ref{fig:psha_map_pga_diff_fsr_poe02} y \ref{fig:psha_map_pga_diff_fsr_poe10}) cuantifican que el efecto máximo se concentra en una franja oriental estrecha y alcanza $2.31\times10^{-2}$~g para PoE 2\% y $1.25\times10^{-2}$~g para PoE 10\% en 50 años. Esta cuantificación es central: el aporte probabilístico de FSR en PGA existe, es positivo y espacialmente coherente, pero no reconfigura la amenaza regional; su efecto queda confinado a un corredor estrecho donde la cercanía a la falla lo hace acumulativamente relevante. Además, el caso BPT muestra una franja más extensa y con mayores valores, consistente con el aumento de contribución relativa cortical observado en desagregación.

\section{Riesgo sísmico}
\label{sec:disc_riesgo}

\subsection{Riesgo por escenarios}
\label{subsec:disc_riesgo_escenarios}


\subsubsection{Nacional versus HAZUS: magnitud de la brecha y lo que implica}
\label{subsubsec:disc_riesgo_vulnerabilidad}

La diferencia entre modelos de vulnerabilidad domina el nivel absoluto de pérdida en los escenarios. A nivel de portafolio, el modelo Nacional entrega pérdidas relativas de 3.70\% (Interplaca), 6.77\% (Intraplaca) y 7.98\% (FSR), mientras que HAZUS entrega 28.67\% (Interplaca), 43.78\% (Intraplaca) y 44.03\% (FSR) (Tabla~\ref{tab:perdida_total_escenarios}). La brecha es de un orden que no puede tratarse como ``variación secundaria'': implica un desplazamiento global del sistema hacia niveles de daño mucho mayores cuando se adopta HAZUS.

Además, el orden entre escenarios cambia con el modelo: en el modelo Nacional, FSR $>$ Intraplaca $>$ Interplaca; en HAZUS, Intraplaca $\approx$ FSR $\gg$ Interplaca. Este punto es crucial porque revela que HAZUS penaliza de manera particularmente fuerte las intensidades características del escenario FSR en el dominio analizado, al punto de compensar su \textit{footprint} más acotado y producir pérdidas agregadas comparables a intraplaca. En otras palabras, el mismo campo de intensidades puede generar narrativas distintas sobre ``qué escenario es más crítico'' dependiendo del modelo de vulnerabilidad, lo que obliga a tratar la elección Nacional/HAZUS como un componente estructural del resultado.

En términos de mecanismo, la propagación desde intensidad a pérdida sigue la lógica de pérdidas esperadas condicionadas a $IM$ (Ecs.~\ref{eq:expected_loss_ratio_from_ds} y \ref{eq:expected_loss_asset}). Por lo tanto, una diferencia de forma en las curvas de vulnerabilidad puede amplificar no linealmente pérdidas en rangos de intensidad donde un escenario concentra su demanda. El hecho de que HAZUS eleve de forma tan marcada el nivel absoluto sugiere que, sin calibración local, su uso debe interpretarse como cota superior metodológica, más que como estimación puntual del nivel esperado de pérdida.

\subsubsection{Sobre el efecto de la materialidad}
\label{subsubsec:disc_riesgo_materialidad}

La desagregación por materialidad (Figura~\ref{fig:perdida_por_materialidad}) muestra que hormigón y albañilería concentran prácticamente la totalidad de la pérdida agregada, mientras madera aporta una fracción marginal. Esta concentración no es sorprendente por sí misma, pero sí es relevante que la partición hormigón/albañilería cambie entre modelos: el modelo Nacional asigna mayor fracción relativa de pérdida a hormigón, mientras HAZUS incrementa el peso relativo de albañilería. Este cambio implica que la selección del modelo de vulnerabilidad no sólo modifica el total, sino también la interpretación de ``qué tipologías explican el riesgo'', y por ende puede orientar de manera distinta la priorización de medidas si el análisis se usa para focalizar intervenciones por tipología.


\subsection{Riesgo probabilístico basado en catálogo de terremotos}
\label{subsec:disc_riesgo_event_based}

\subsubsection{AAL y AALR: señal espacial y coherencia con la dominancia de subducción}
\label{subsubsec:disc_aal_aalr}

Los mapas sin FSR (Figuras~\ref{fig:aal_mapas_sin_fsr} y \ref{fig:aalr_mapas_sin_fsr}) muestran una separación clara entre modelos: HAZUS presenta valores más altos y mayor extensión de manzanas en categorías intermedias-altas, mientras el modelo Nacional concentra el dominio en rangos bajos. Al incluir FSR (Figuras~\ref{fig:aal_mapas_con_fsr} y \ref{fig:aalr_mapas_con_fsr}), la estructura espacial global se mantiene, lo que es coherente con el diagnóstico PSHA de dominancia de subducción: una fuente adicional cortical modifica la amenaza y el riesgo de forma localizada, pero no redefine el patrón global.

A nivel de portafolio, los cambios \pFifty{} por incluir FSR son modestos (Tablas~\ref{tab:resumen_portafolio_aal_aalr} y \ref{tab:diferencia_portafolio_fsr}). En el modelo Nacional, AAL \pFifty{} pasa de $9.14\times10^{8}$ a $9.31\times10^{8}$~CLP/año, con $\Delta$AAL de $1.70\times10^{7}$~CLP/año (1.86\%) y $\Delta$AALR de $2.0\times10^{-2}$\textperthousand\ (2.27\%). En HAZUS, AAL \pFifty{} pasa de $7.70\times10^{9}$ a $7.76\times10^{9}$~CLP/año, con $\Delta$AAL de $6.00\times10^{7}$~CLP/año (0.78\%) y $\Delta$AALR de $5.0\times10^{-2}$\textperthousand\ (0.67\%). La lectura incisiva es que el efecto FSR en la esperanza anual agregada existe y es positivo, pero su magnitud queda en el orden de pocos puntos porcentuales, alineado con que la amenaza es controlada por subducción.

\subsubsection{Efecto de FSR en AALR: por qué es más visible localmente}
\label{subsubsec:disc_efecto_fsr_aalr}

El mapa de incremento relativo de AALR por incluir FSR (Figura~\ref{fig:aalr_efecto_fsr_pct}) muestra aumentos positivos en todo el dominio y máximos locales del orden de 2.1\%.


\subsubsection{\texorpdfstring{Curvas OEP y tasas $\lambda(L)$: dónde se manifiesta el aporte cortical}{Curvas OEP y tasas lambda(L): donde se manifiesta el aporte cortical}}
\label{subsubsec:disc_oep_lambda}

Las curvas OEP relativas (Figura~\ref{fig:oep_relativo_con_sin_fsr}) muestran una separación consistente entre modelos y entre subconjuntos del inventario: HAZUS predice excedencias relativas mayores a lo largo del rango de periodos de retorno, y el subconjunto sin FSR reduce sistemáticamente los niveles respecto a considerar todas las fuentes. En paralelo, las tasas $\lambda(L)$ (Figura~\ref{fig:lambda_comparison_con_sin_fsr}) muestran que, para un umbral relativo dado, HAZUS mantiene tasas mayores y la inclusión de FSR desplaza la tasa hacia arriba. El efecto marginal $\Delta\lambda(L)$ es positivo en todo el rango y decrece hacia cero al aumentar el umbral, lo que sugiere un punto clave: la FSR contribuye proporcionalmente más a excedencias frecuentes de pérdidas bajas-medias que a la cola extrema. Esto es consistente con la combinación de una fuente localizada que incrementa el número de eventos con pérdidas moderadas en las inmediaciones de la falla, mientras que las pérdidas más extremas están dominadas por eventos de subducción capaces de generar pérdidas grandes a escala regional.


\section{Conclusiones}
\label{sec:conclusiones}

\subsection{Amenaza}
\label{subsec:conc_amenaza}

\subsubsection{FSR domina la amenaza determinista local, subducción domina la amenaza probabilística regional}
\label{subsubsec:conc_amenaza_fsr_subd}

La amenaza determinista muestra que la FSR produce máximos y gradientes longitudinales fuertes hacia el margen oriental (Figuras~\ref{fig:dsha_map_pga_all} y \ref{fig:dsha_map_sa10_all}), con incertidumbre y sensibilidad geométrica que se manifiestan localmente (Figuras~\ref{fig:dsha_map_pga_fsr_percentiles} y \ref{fig:dsha_map_pga_fsr_dip_traza}). En cambio, la amenaza probabilística en el rango de PoE reportado está dominada por subducción, con FSR como aporte menor en curvas y desagregación, salvo en sitios específicos donde su contribución aumenta bajo BPT (Figuras~\ref{fig:psha_curvas_pga_poisson}--\ref{fig:psha_disagg_sa10_bpt}). Esta combinación no es contradictoria: describe dos escalas distintas. El DSHA identifica dónde la demanda máxima plausible se concentra; el PSHA muestra que, al promediar por recurrencia, subducción controla la mayor parte de la excedencia en el horizonte analizado.

\subsubsection{El modelo temporal (BPT) aumenta la relevancia relativa de FSR sin alterar la dominancia interplaca e intraplaca}
\label{subsubsec:conc_amenaza_bpt}

La dependencia temporal realza la contribución relativa de FSR en sitios más sensibles, especialmente en $SA(1.0)$ en Sitios 1 y 2, donde alcanza 13\% en PoE 2\% en 50 años (Figura~\ref{fig:psha_disagg_sa10_bpt}). Sin embargo, la interplaca permanece como aporte mayoritario en $SA(1.0)$ y la estructura espacial regional se mantiene, indicando que el modelo temporal ajusta la contribución relativa sin cambiar la dominancia.

\subsection{Riesgo}
\label{subsec:conc_riesgo}

\subsubsection{La vulnerabilidad controla el nivel absoluto y condiciona la interpretación tipológica}
\label{subsubsec:conc_riesgo_vuln}

El contraste Nacional versus HAZUS es el factor más influyente en el nivel absoluto de pérdida, tanto en riesgo por escenarios como en métricas anualizadas y de excedencia. La brecha observada en escenarios (Tabla~\ref{tab:perdida_total_escenarios}) y su persistencia en AAL/AALR y $\lambda(L)$ indican que la incertidumbre epistemológica dominante proviene de la vulnerabilidad, no de la amenaza. Además, la partición por materialidad cambia entre modelos, afectando la narrativa de cuáles tipologías explican el riesgo agregado (Figura~\ref{fig:perdida_por_materialidad}).

\subsubsection{El efecto de FSR en portafolio es pequeño, pero su señal es consistente y localmente relevante}
\label{subsubsec:conc_riesgo_fsr}

En el enfoque \textit{event-based}, la inclusión de FSR incrementa AAL y AALR de portafolio en pocos puntos porcentuales (Tabla~\ref{tab:diferencia_portafolio_fsr}), coherente con la dominancia de subducción. Sin embargo, el incremento local en AALR es espacialmente coherente y alcanza del orden de 2.1\% en manzanas (Figura~\ref{fig:aalr_efecto_fsr_pct}), mostrando que el efecto cortical se manifiesta donde la proximidad a la traza lo concentra. El comportamiento de $\Delta\lambda(L)$ decreciente sugiere que la FSR tiene más peso relativo en pérdidas frecuentes bajas-medias que en la cola extrema, dominada por subducción (Figura~\ref{fig:delta_lambda_con_menos_sin_fsr}).


\section{Limitaciones y proyecciones}
\label{sec:limitaciones_proyecciones}

% \subsection{Limitaciones}
% \label{subsec:limitaciones}

\subsection{Interpretación del nivel absoluto bajo HAZUS}
\label{subsubsec:lim_hazus}

Las pérdidas relativas de HAZUS en escenarios (28\%--44\%) son sustantivamente mayores que las del modelo Nacional (3.7\%--8.0\%) para el mismo dominio e inventario (Tabla~\ref{tab:perdida_total_escenarios}). Esta magnitud sugiere que HAZUS debe interpretarse con cautela si el objetivo es estimación puntual del nivel absoluto, y es más defendible como referencia comparativa o cota metodológica, en ausencia de calibración local. La consecuencia práctica es que decisiones basadas en umbrales absolutos de pérdida pueden cambiar de forma drástica según el modelo.

\subsection{Dependencia temporal en FSR}
\label{subsubsec:lim_bpt}

El aumento de contribución FSR bajo BPT depende de la parametrización temporal y, por lo tanto, su interpretación debe leerse como sensibilidad epistemológica del modelo cortical. El resultado es robusto en dirección (aumento relativo en sitios cercanos), pero su magnitud específica requiere ser contrastada con evidencia geológica y con análisis de sensibilidad de parámetros si se busca una recomendación más prescriptiva.

% \subsection{Proyecciones}
% \label{subsec:proyecciones}

\subsection{Calibración local de vulnerabilidad y validación}
\label{subsubsec:proj_vuln}

La prioridad más directa para reducir incertidumbre en el nivel absoluto es avanzar en calibración local de vulnerabilidad, especialmente en albañilería y hormigón, dado que la brecha Nacional vs.\ HAZUS domina resultados anualizados y por escenarios. Sin esta calibración, el análisis entrega con mayor robustez la localización relativa del riesgo que su magnitud absoluta.
